% GENERATED FILE. Edit canonical lessons, then run npm run generate:books.
% canonical-source-hash: fnv1a64:d0f6c35c3b56342c
% canonical-lessons: UR-C253-chhipkali, UR-C253-kachhua, UR-C253-magarmachh, UR-C253-zarafa, UR-C253-makri

\chapter{Lizard, Tortoise, Crocodile, Giraffe, Spider}
\label{ch:ur-lizard-tortoise-crocodile-giraffe-spider}
\hlchaptermodality{253}

\begin{chapteropening}
\textbf{By the end of this chapter:} \emph{I can say a lizard, a tortoise, a crocodile, a giraffe and a spider.}

Complete the last lesson of chapter 253: I can say a lizard, a tortoise, a crocodile, a giraffe and a spider.
\end{chapteropening}

\section[\texorpdfstring{\ur{چھپکلی} (chhipkalī)}{چھپکلی (chhipkalī)}]{\ur{چھپکلی} (chhipkalī) — a lizard}

\label{lesson:UR-C253-chhipkali}



\begin{quote}
Before the new one: say the Urdu for a coconut, then the Urdu for a cashew.
\end{quote}

\subsection*{You'll want to know: \ur{چھپکلی}}
\textbf{\ur{چھپکلی}} — \emph{chhipkalī} — \textquotedblleft{}a lizard\textquotedblright{}.

\begin{cousinweb}
CDIAL lists Hindi \emph{chipkalī}, \textquotedblleft{}lizard\textquotedblright{}, under Sanskrit \emph{śepyā}, for which it reconstructs a sense \textquotedblleft{}tail\textquotedblright{} (CDIAL 12607). How the word came to name the lizard is not explained.
\end{cousinweb}

\begin{grammarlens}[title={What you've built}]
One more word for this chapter.
\end{grammarlens}

\subsection*{Your turn}
\begin{itemize}
\raggedright
  \item \emph{Say it:} \emph{chhipkalī}
  \item \emph{Say it:} \emph{chhipkalī}, once more
  \item \emph{Say it:} say \emph{chhipkalī}
  \item \emph{Recall:} say the Urdu for a coconut, then the Urdu for a cashew, then say \emph{chhipkalī} again
\end{itemize}

\begin{tcolorbox}[breakable,colback=teal!4,colframe=teal!35!black,title={Before you move on}]
What does \ur{چھپکلی} mean? (\textquotedblleft{}A lizard\textquotedblright{}.) Say it once more.
\end{tcolorbox}
\section[\texorpdfstring{\ur{کچھوا} (kachhuā)}{کچھوا (kachhuā)}]{\ur{کچھوا} (kachhuā) — a tortoise}

\label{lesson:UR-C253-kachhua}



\begin{quote}
Before the new one: say the Urdu for a cashew, then the Urdu for a lizard.
\end{quote}

\subsection*{You'll want to know: \ur{کچھوا}}
\textbf{\ur{کچھوا}} — \emph{kachhuā} — \textquotedblleft{}a tortoise\textquotedblright{}.

\begin{cousinweb}
\textbf{\ur{کچھوا}} (\emph{kachhuā}) is a tadbhava, a word inherited from Sanskrit and reshaped by sound change. It comes from \emph{kacchapa}, \textquotedblleft{}turtle, tortoise\textquotedblright{} (CDIAL 2619).
\end{cousinweb}

\begin{grammarlens}[title={What you've built}]
One more word for this chapter.
\end{grammarlens}

\subsection*{Your turn}
\begin{itemize}
\raggedright
  \item \emph{Say it:} \emph{kachhuā}
  \item \emph{Say it:} \emph{kachhuā}, once more
  \item \emph{Say it:} say \emph{kachhuā}
  \item \emph{Recall:} say the Urdu for a cashew, then the Urdu for a lizard, then say \emph{kachhuā} again
\end{itemize}

\begin{tcolorbox}[breakable,colback=teal!4,colframe=teal!35!black,title={Before you move on}]
What does \ur{کچھوا} mean? (\textquotedblleft{}A tortoise\textquotedblright{}.) Say it once more.
\end{tcolorbox}
\section[\texorpdfstring{\ur{مگرمچھ} (magarmachh)}{مگرمچھ (magarmachh)}]{\ur{مگرمچھ} (magarmachh) — a crocodile}

\label{lesson:UR-C253-magarmachh}



\begin{quote}
Before the new one: say the Urdu for a lizard, then the Urdu for a tortoise.
\end{quote}

\subsection*{You'll want to know: \ur{مگرمچھ}}
\textbf{\ur{مگرمچھ}} — \emph{magarmachh} — \textquotedblleft{}a crocodile\textquotedblright{}.

\begin{cousinweb}
\textbf{\ur{مگرمچھ}} (\emph{magarmachh}) is \textquotedblleft{}crocodile-fish\textquotedblright{}. \emph{magar} goes back to Sanskrit \emph{makara}, \textquotedblleft{}crocodile\textquotedblright{}, a word CDIAL records without a Hindi form (CDIAL 9692), and \emph{machh}, \textquotedblleft{}fish\textquotedblright{}, to Sanskrit \emph{matsya} (CDIAL 9758).
\end{cousinweb}

\begin{grammarlens}[title={What you've built}]
One more word for this chapter.
\end{grammarlens}

\subsection*{Your turn}
\begin{itemize}
\raggedright
  \item \emph{Say it:} \emph{magarmachh}
  \item \emph{Say it:} \emph{magarmachh}, once more
  \item \emph{Say it:} say \emph{magarmachh}
  \item \emph{Recall:} say the Urdu for a lizard, then the Urdu for a tortoise, then say \emph{magarmachh} again
\end{itemize}

\begin{tcolorbox}[breakable,colback=teal!4,colframe=teal!35!black,title={Before you move on}]
What does \ur{مگرمچھ} mean? (\textquotedblleft{}A crocodile\textquotedblright{}.) Say it once more.
\end{tcolorbox}
\section[\texorpdfstring{\ur{زرافہ} (zarāfā)}{زرافہ (zarāfā)}]{\ur{زرافہ} (zarāfā) — a giraffe}

\label{lesson:UR-C253-zarafa}



\begin{quote}
Before the new one: say the Urdu for a tortoise, then the Urdu for a crocodile.
\end{quote}

\subsection*{You'll want to know: \ur{زرافہ}}
\textbf{\ur{زرافہ}} — \emph{zarāfā} — \textquotedblleft{}a giraffe\textquotedblright{}.

\begin{cousinweb}
\textbf{\ur{زرافہ}} (\emph{zarāfā}) is a loan from Arabic \emph{zarāfa}, \textquotedblleft{}giraffe\textquotedblright{}, which reached Urdu through Persian. Its Arabic root is \emph{z-r-f} (Lane). English \emph{giraffe} comes from the same Arabic word.
\end{cousinweb}

\begin{grammarlens}[title={What you've built}]
One more word for this chapter.
\end{grammarlens}

\subsection*{Your turn}
\begin{itemize}
\raggedright
  \item \emph{Say it:} \emph{zarāfā}
  \item \emph{Say it:} \emph{zarāfā}, once more
  \item \emph{Say it:} say \emph{zarāfā}
  \item \emph{Recall:} say the Urdu for a tortoise, then the Urdu for a crocodile, then say \emph{zarāfā} again
\end{itemize}

\begin{tcolorbox}[breakable,colback=teal!4,colframe=teal!35!black,title={Before you move on}]
What does \ur{زرافہ} mean? (\textquotedblleft{}A giraffe\textquotedblright{}.) Say it once more.
\end{tcolorbox}
\section[\texorpdfstring{\ur{مکڑی} (makṛī)}{مکڑی (makṛī)}]{\ur{مکڑی} (makṛī) — a spider}

\label{lesson:UR-C253-makri}



\begin{quote}
Before the new one: say the Urdu for a crocodile, then the Urdu for a giraffe.
\end{quote}

\subsection*{You'll want to know: \ur{مکڑی}}
\textbf{\ur{مکڑی}} — \emph{makṛī} — \textquotedblleft{}a spider\textquotedblright{}.

\begin{cousinweb}
\textbf{\ur{مکڑی}} (\emph{makṛī}) is a tadbhava, a word inherited from Sanskrit and reshaped by sound change. It comes from \emph{markaṭa}, \textquotedblleft{}spider\textquotedblright{} (CDIAL 9883). The Sanskrit word is known only from the old word-lists.
\end{cousinweb}

\begin{grammarlens}[title={What you've built}]
One more word for this chapter.
\end{grammarlens}

\subsection*{Your turn}
\begin{itemize}
\raggedright
  \item \emph{Say it:} \emph{makṛī}
  \item \emph{Say it:} \emph{makṛī}, once more
  \item \emph{Say it:} say \emph{makṛī}
  \item \emph{Recall:} say the Urdu for a crocodile, then the Urdu for a giraffe, then say \emph{makṛī} again
\end{itemize}

\begin{tcolorbox}[breakable,colback=teal!4,colframe=teal!35!black,title={Before you move on}]
What does \ur{مکڑی} mean? (\textquotedblleft{}A spider\textquotedblright{}.) Say it once more.
\end{tcolorbox}
